%!TEX root = ../QuantumNoiseAndDecoherence.tex
% ──────────────────────────────────────────────────────────────
%  Lecture 2 — Entanglement, Dynamics, and Quantum Channels
% ──────────────────────────────────────────────────────────────

% Continues Section 1 (Recapitulation of QM) from Lecture 1.

\subsection{Entanglement versus classical correlations}%
\label{sec:entanglement-vs-correlations}

We saw in Lecture~1 that the Schmidt decomposition tells us when a
\emph{pure} state of a composite system is entangled: it is entangled if
and only if its Schmidt rank is greater than one.  But for
\emph{mixed} states, the situation is more subtle.

Consider the two-qubit density operator
\begin{equation}\label{eq:classical-corr}
  \rho = p\,\ketbra{00}{00} + (1-p)\,\ketbra{11}{11}\,,
  \qquad p \in (0,1)\,.
\end{equation}
This is \emph{not} a product state: $\rho \neq \rho_A \tp \rho_B$
for any pair of single-qubit states $\rho_A$, $\rho_B$.  Yet it is
manifestly \emph{not entangled}---it can be prepared by a purely
classical procedure:

\begin{intuition}[Correlated coins]
  Flip a (biased) coin.  If heads (probability~$p$), prepare both
  qubits in $\ket{0}$; if tails (probability~$1-p$), prepare both
  in~$\ket{1}$.  Hand one qubit to Alice and the other to Bob.  The
  joint state is~\eqref{eq:classical-corr}---correlated, but only
  classically so.
\end{intuition}

A state is called \textbf{separable} (or \emph{classically correlated})
if it can be written as a convex combination of product states:
\begin{equation}\label{eq:separable}
  \rho = \sum_k \lambda_k\, \rho_A^{(k)} \tp \rho_B^{(k)}\,,\qquad
  \lambda_k \geq 0,\quad \sum_k \lambda_k = 1\,.
\end{equation}
A state that is \emph{not} separable is \textbf{entangled}.

\begin{remark}
  For pure states, ``separable'' reduces to ``product,'' and
  entanglement can be efficiently decided via the Schmidt
  decomposition.  For mixed states, deciding whether a given~$\rho$ is
  separable or entangled is, in general, a computationally hard problem
  and an active area of research in quantum information theory.
\end{remark}

\subsubsection{Entanglement and quantum noise}

There is a deep connection between entanglement and quantum noise that
motivates this entire course.  In classical probability, a pure state
of a composite system always assigns definite (pure) states to each
subsystem.  In quantum mechanics, this is spectacularly \emph{false}:

\begin{keyresult}[Entanglement $\Leftrightarrow$ noise on subsystems]
  A composite quantum system can be in a pure state (zero temperature,
  unitary evolution, no noise whatsoever), yet its subsystems are in
  \emph{mixed} states---they look noisy and are decohering rapidly.
  This happens precisely when the global state is entangled.
\end{keyresult}

Recall from the Schmidt decomposition (Proposition~\ref{prop:schmidt})
that for any pure state $\ket{\psi} \in \hilbA \tp \hilbB$ with
Schmidt form $\ket{\psi} = \sum_j \sqrt{p_j}\,\ket{g_j}\tp\ket{f_j}$,
the reduced state of subsystem~$A$ is
\begin{equation}\label{eq:reduced-from-schmidt}
  \rho_A = \tr_B(\ketbra{\psi}{\psi})
         = \sum_j p_j\,\ketbra{g_j}{g_j}\,.
\end{equation}
This is pure if and only if exactly one~$p_j$ is nonzero---i.e.\ when
$\ket{\psi}$ is a product state.  Otherwise, $\rho_A$ is mixed:
the subsystem has acquired ``noise'' entirely from entanglement with
the rest of the system.

This is, in some sense, the central theme of this course: noise comes
from entanglement with an environment, and quantum computation exploits
entanglement within a system.  Making a quantum computer work requires
fighting entanglement with entanglement.

% ──────────────────────────────────────────────────────────────
\subsection{Dynamics of quantum systems}\label{sec:dynamics}

A good physical theory has both \emph{kinematics} (how to represent
states and observables) and \emph{dynamics} (how to represent the
passage of time).  The fundamental dynamical question is:

\begin{quote}
  \emph{If I know how to represent a measurement at time $t=0$, how
  do I represent the same measurement at time~$t$?}
\end{quote}

We divide a generic experiment into three stages:

\begin{center}
\begin{tikzpicture}[
  block/.style={ellipse, draw=cgblue, fill=cgblue!10,
    minimum height=2em, minimum width=3cm,
    font=\sf\small, text=spacecadet, inner sep=5pt},
  arr/.style={-{Stealth[length=5pt]}, thick, cgblue}
]
  \node[block] (prep) {Preparation};
  \node[block, right=1.5cm of prep] (evol) {Evolution};
  \node[block, right=1.5cm of evol] (meas) {Measurement};
  \draw[arr] (prep) -- (evol);
  \draw[arr] (evol) -- (meas);
\end{tikzpicture}
\end{center}

\subsubsection{Heisenberg and Schr\"odinger pictures}%
\label{sec:pictures}

There are two complementary ways to absorb the evolution back into the
kinematics:

\begin{itemize}
  \item \textbf{Schr\"odinger picture:} the evolution produces a new
    effective state $\omega'$ from the old state~$\omega$; observables
    are unchanged.
  \item \textbf{Heisenberg picture:} the evolution produces a new
    effective observable~$A'$ from the original~$A$; the state is
    unchanged.
\end{itemize}
Both pictures describe the same physics.  Their compatibility is
expressed by requiring that for every state~$\rho$ and
observable~$A$,
\begin{equation}\label{eq:pictures-compat}
  \tr\bigl(\rho\;\mathcal{U}(A)\bigr)
    = \tr\bigl(\mathcal{U}_*({\rho})\;A\bigr)\,,
\end{equation}
where $\mathcal{U}$ is the Heisenberg-picture evolution and
$\mathcal{U}_*$ is its Schr\"odinger-picture dual.

\begin{remark}
  There are also \emph{interaction pictures} in which part of the
  evolution is applied to the state and part to the observable.  These
  are useful for perturbation theory but will not play a role in the
  first part of this course.
\end{remark}

\subsubsection{Requirements for evolution maps}%
\label{sec:evolution-requirements}

What are the minimal requirements for a map to qualify as a physical
evolution?  Working in the Heisenberg picture, an evolution
$\mathcal{U}\colon \alg \to \alg$ must satisfy:

\begin{enumerate}[(i)]
  \item \textbf{Linearity:} $\mathcal{U}$ is a linear map.  This
    follows from the requirement that randomising over preparations
    before or after the evolution gives the same result.
  \item \textbf{Probability preservation (unitality):}
    $\mathcal{U}(\One) = \One$.  If the effect ``do nothing'' is
    evolved, it must remain ``do nothing''---otherwise probability is
    created or destroyed.
  \item \textbf{Positivity:} if $E \geq 0$ then
    $\mathcal{U}(E) \geq 0$.  Legal effects must be mapped to legal
    effects.
\end{enumerate}

\subsubsection{Closed systems: unitary evolution}%
\label{sec:unitary-evol}

A \emph{closed} (or \emph{reversible}) evolution is one that, in
addition to the above, is invertible.  For finite-dimensional systems,
closed evolutions are exactly those representable in terms of a
unitary operator $U$:
\begin{equation}\label{eq:unitary-heisenberg}
  \eqbox{\mathcal{U}(A) = U^\adj A\, U}\,,\qquad U^\adj U = U U^\adj = \One\,.
\end{equation}
In the Schr\"odinger picture, the corresponding map is
\begin{equation}\label{eq:unitary-schrodinger}
  \mathcal{U}_*(\rho) = U\,\rho\, U^\adj\,.
\end{equation}

\begin{remark}
  In infinite dimensions, there exist reversible dynamics that cannot be
  represented as conjugation by a unitary on a Hilbert space.  The
  first place this arises is for infinite tensor products of qubits
  (spin chains).  We will not encounter this subtlety in this course.
\end{remark}

\subsubsection{The Schr\"odinger equation}%
\label{sec:schrodinger-eq}

If the unitary $U(t)$ depends on a time parameter~$t$ in a
\emph{time-homogeneous} fashion, meaning
\begin{equation}\label{eq:time-homogeneous}
  U(t+s) = U(t)\,U(s)\,,\qquad U(0) = \One\,,
\end{equation}
then $U(t)$ forms a one-parameter group and satisfies the
\textbf{Schr\"odinger equation}:
\begin{equation}\label{eq:schrodinger}
  \eqbox{i\,\frac{\dd}{\dd t}\,U(t) = H\, U(t)}\,,
\end{equation}
where $H = H^\adj$ is the \emph{Hamiltonian}, a self-adjoint operator
generating the evolution.  The solution is
$U(t) = e^{-iHt}$.

% ──────────────────────────────────────────────────────────────
\subsection{Open quantum systems}\label{sec:open-systems}

Is the Schr\"odinger equation the whole story?  \emph{Yes and no.}
The aim of this course is to study the dynamics of \emph{composite}
quantum systems, and this leads us far beyond simple unitary evolution.

\begin{center}
\begin{tikzpicture}[
  sysbox/.style={rectangle, rounded corners=3pt, draw=cgblue,
    fill=cgblue!8, minimum height=2.5em, minimum width=2cm,
    font=\sf\small, text=spacecadet, inner sep=6pt},
  envbox/.style={rectangle, rounded corners=3pt, draw=cgblue!50,
    fill=isabelline, minimum height=2.5em, minimum width=2cm,
    font=\sf\small, text=spacecadet!70, inner sep=6pt},
  bigbox/.style={rectangle, rounded corners=5pt, draw=cgblue,
    dashed, inner sep=10pt},
  arr/.style={-{Stealth[length=5pt]}, thick, cgblue}
]
  \node[sysbox] (A) {System $A$};
  \node[envbox, right=1.5cm of A] (E) {Environment $E$};
  \node[bigbox, fit=(A)(E), label={[font=\sf\small,
    text=spacecadet]above:Universe $AE$}] (AE) {};
  % Interaction arrow
  \draw[arr, <->] (A) -- (E)
    node[midway, above, font=\footnotesize]{$U(t)$};
\end{tikzpicture}
\end{center}

The picture is as follows.  Our system~$A$ (an atom in a trap, a mode of
light, a qubit) is embedded in a larger system~$AE$ that includes the
environment~$E$.  The \emph{universe}~$AE$ undergoes reversible
(unitary) evolution~$U(t)$.  Our task: understand what this evolution
looks like when \emph{restricted to~$A$}.

\begin{keyresult}[Irreversibility from ignoring subsystems]
  Even though the global evolution $U(t)$ is unitary and reversible,
  the effective evolution on subsystem~$A$ alone is, in general,
  \emph{irreversible}.  Irreversibility arises from tracing out
  (ignoring) the environment.
\end{keyresult}

This is precisely how thermodynamics works: a system interacts with a
reservoir; the joint evolution is unitary; but when we focus on the
system alone, we see irreversible behaviour---dissipation, decoherence,
thermalisation.

\begin{intuition}[The Lindblad equation]
  Is there a ``Schr\"odinger equation just for~$A$''?  Sometimes
  yes---under a Markov approximation, the effective evolution of~$A$
  obeys a \textbf{Lindblad equation} (or master equation).  Deriving
  and understanding this equation will be one of the main results of
  this course.
\end{intuition}

It used to be that open quantum systems were treated as an annoying
detail to be left to a PhD student.  This is no longer the case:
dissipative dynamics, arising from interactions with a noisy
environment, can actually be \emph{useful}---they can prepare desired
states and perform computations.

% ──────────────────────────────────────────────────────────────
\subsection{Quantum channels}\label{sec:channels}

We now axiomatically characterise the most general evolution in quantum
mechanics: the \emph{quantum channel}.  Two features distinguish
channels from the reversible evolutions of Section~\ref{sec:unitary-evol}:
\begin{enumerate}[(a)]
  \item The input and output systems may \emph{differ}: a channel can
    map observables on a system~$B$ to observables on a system~$A$,
    where $A$ and~$B$ need not be the same.
  \item In addition to positivity, we require \emph{complete
    positivity}.
\end{enumerate}

\begin{definition}[Quantum channel]\label{def:channel}
  A \emph{quantum channel} (in the Heisenberg picture) is a linear map
  \[
    \chan\colon \algB \to \alg
  \]
  satisfying:
  \begin{enumerate}[(1)]
    \item \textbf{Unitality (probability preservation):}
      $\chan(\One_B) = \One_A$.
    \item \textbf{Complete positivity:} for every $n \in \N$, the
      extended map
      $\chan \tp \id_n\colon \algB \tp \Mn{n} \to \alg \tp \Mn{n}$
      is positive:
      \begin{equation}\label{eq:cp}
        X \geq 0 \;\Longrightarrow\;
        (\chan \tp \id_n)(X) \geq 0\,.
      \end{equation}
  \end{enumerate}
\end{definition}

\subsubsection{Why does $\chan$ map $\algB \to \alg$?}

In the Heisenberg picture, the channel acts on observables in the
\emph{reverse} direction to the flow of time:

\begin{center}
\begin{tikzpicture}[
  block/.style={rectangle, rounded corners=3pt, draw=cgblue,
    fill=cgblue!8, minimum height=2em, minimum width=2cm,
    font=\sf\small, text=spacecadet, inner sep=5pt},
  arr/.style={-{Stealth[length=5pt]}, thick, cgblue}
]
  \node[block] (prep) {$\omega$ on $A$};
  \node[block, right=2.5cm of prep] (meas) {Effect $F$ on $B$};
  \draw[arr] (prep) -- (meas)
    node[midway, above, font=\footnotesize]{time};
  \draw[arr, munsell, bend left=30] (meas) to
    node[midway, below, font=\footnotesize, text=munsell]
      {$\chan(F) \in \alg$} (prep);
\end{tikzpicture}
\end{center}

If we prepare state~$\omega$ on~$A$ and, after the evolution, measure
effect~$F$ on~$B$, then $\chan(F)$ is the effective observable on~$A$
that gives the same statistics: $\omega(\chan(F)) = p(F)$.

\subsubsection{Why complete positivity?}\label{sec:why-cp}

Ordinary positivity ($E \geq 0 \Rightarrow \chan(E) \geq 0$) is not
enough.  Suppose we perform our evolution on system~$A$ while an
ancillary system~$N$ sits idle in a neighbouring laboratory.  The joint
evolution is $\chan \tp \id_N$.  If we then perform a \emph{correlated}
measurement on the joint system~$AN$, we need the result to be a legal
(positive) effect.  This requires $\chan \tp \id_N$ to be positive for
\emph{every} ancillary system---which is precisely complete positivity.

\begin{example}[The transpose map]\label{ex:transpose}
  The map $T\colon A \mapsto A^\top$ (transpose in a fixed basis) is:
  \begin{itemize}
    \item linear,
    \item unital: $T(\One) = \One^\top = \One$,
    \item positive: $A \geq 0 \Rightarrow A^\top \geq 0$ (since
      eigenvalues are preserved by transposition).
  \end{itemize}
  However, $T$ is \emph{not} completely positive.  The extended map
  $T \tp \id_n$ fails to be positive: it maps certain positive
  elements of $\Mn{n} \tp \Mn{n}$ to operators with negative
  eigenvalues.  Therefore, the transpose is \textbf{not} a valid
  quantum channel.
\end{example}

\begin{exercise}\label{ex:transpose-not-cp}
  Let $\ket{\Phi^+} = \tfrac{1}{\sqrt{2}}(\ket{00} + \ket{11})$ be a
  maximally entangled state of two qubits.  Show that
  $(T \tp \id)(\ketbra{\Phi^+}{\Phi^+})$ has a negative eigenvalue,
  thereby proving that the transpose map is not completely positive.

  \emph{Hint:} Compute the matrix explicitly and find its eigenvalues.
\end{exercise}
